\documentclass[11pt]{article}
\usepackage{icmlko}
\begin{document}

\title{판의 단위는 도메인 전체였다 --- 어제 만든 걱정 하나를 지운다}
\author{월드모델 실험실}
\date{노트 469 · 2026년 8월 2일}
\maketitle

\begin{abstract}
노트 468 이 \emph{``집 밖 효과의 부호가 채점 범위에 딸렸다''}를 보이고
남긴 것 중 제일 큰 것은 이것이었다 --- \emph{``판은 도메인 \textbf{안}
순위를 도메인마다 재서 가중평균한다. 즉 갈래로 가른 쪽에 가깝다. 집 밖
수치를 판과 견줄 때 이 차이를 안 적고 있었다.''}

그게 맞다면 심각하다. 노트 $459 \sim 460$ 의 \textbf{환율}은 판 손해를
집 밖 이득으로 나눈 값이고, 노트 457 의 깊이 $12$ 채택이 환율 $1.95$
위에 선다. 범위가 어긋났다면 그 나눗셈이 무의미하다.

그래서 범위를 맞춰 다시 쟀고, \textbf{정말로 뒤집혔다.}

\begin{center}
\begin{tabular}{llrr}
\hline
짝 & 범위 & 깊이 $6 \to 12$ & 95\% 구간 \\
\hline
KR 만화 & 수집 전체 & $+0.0048$ & $[+0.0003,+0.0093]$ \\
KR 만화 & 갈래 안 & $\mathbf{-0.0090}$ & $[-0.0114,-0.0065]$ \\
비게임 앱 & 수집 전체 & $+0.0048$ & $[+0.0012,+0.0083]$ \\
비게임 앱 & 갈래 안 & $+0.0045$ & $[-0.0000,+0.0090]$ \\
\hline
\end{tabular}
\end{center}

\textbf{그런데 전제가 틀렸다.} 코드를 열어 확인했다.
\end{abstract}

\begin{figure}[t]
\centering
\includegraphics[width=\textwidth]{figs/f1.pdf}
\caption{\textbf{가} --- 깊이도 범위에 따라 부호가 바뀐다(노트 468 의
현상이 가림에만 있는 게 아니다). \textbf{나} --- 그러나 판의 단위는
도메인 \emph{전체}였다.}
\end{figure}

\section{\texttt{\_score\_one} 이 답이었다}

\begin{center}
\texttt{return float(spearman(p[ok], y[ok]))}
\end{center}

판은 도메인마다 \textbf{그 도메인의 유보 행 전부}에 스피어만을 하나
매기고, 그것을 채점 표본 수로 가중평균한다. \emph{갈래로 안 가른다.}

내가 논문 354 에서 ``도메인 \textbf{안} 순위''라고 쓴 것은 맞는 말인데,
그 다음 문장에서 그것을 \emph{``갈래로 가른 쪽''}으로 읽었다.
\textbf{``도메인 안''과 ``갈래 안''은 다르다.} 도메인 안이란 곧 그 도메인
\emph{전체}이고, 집 밖 짝에서 그것에 대응하는 것은 \textbf{수집 전체}다.

\textbf{범위는 처음부터 맞아 있었다.} 노트 $459 \sim 460$ 의 환율도,
노트 457 의 깊이 $12$ 채택도 어긋난 나눗셈 위에 서 있지 않다.
\textbf{걱정을 지운다.}

\section{그래도 남는 것 --- 범위 딸림은 가림만의 일이 아니다}

전제가 틀렸다고 측정이 버려지지는 않는다. 위 표는 \textbf{깊이도} 범위에
따라 부호를 바꾼다는 것을 보여 준다 --- KR 에서 수집 전체 $+0.0048$,
갈래 안 $-0.0090$, 두 구간 모두 $0$ 을 안 문다.

노트 468 은 이 현상을 \texttt{entry\_friction} 가림 하나에서 봤다. 이제
\textbf{용량 손잡이에서도} 같은 일이 일어난다는 것을 안다. \emph{범위
딸림은 어떤 특정 변경의 성질이 아니라 이 자료의 성질이다.}

노트 468 의 조항은 그대로 선다 --- \textbf{집 밖 효과를 적을 때 채점
범위를 같이 적는다.} 다만 한 줄을 고쳐 붙인다.

\begin{quote}
\textbf{정정} --- 지금까지 적은 집 밖 수치는 전부 \textbf{수집 전체}
범위이고, 판의 도메인 점수와 \textbf{같은 범위}다. 견줘도 된다. 범위를
적는 이유는 어긋나서가 아니라 \emph{다른 범위의 수를 나중에 섞지
않으려고}다.
\end{quote}

\section{두 번째로 같은 실수를 했다}

이 사이클에서 내 설명이 틀린 게 세 번째다 --- 노트 450(더 채워서 이겼다),
노트 467(쓰임 규칙), 그리고 이번. 앞의 둘은 \textbf{측정}으로 잡았고
이번은 \textbf{코드를 읽어서} 잡았다.

\textbf{다른 점이 하나 있다.} 앞의 둘은 틀린 설명이 다음 실험을 낳았고
그 실험이 스스로를 죽였다. 이번은 틀린 전제가 \emph{채택 하나를 되돌릴}
뻔했다 --- 노트 469 의 사전 등록대로면 깊이 $12$ 를 되돌렸을 것이다.
\textbf{되돌리기 전에 코드를 열어 본 것이 유일하게 옳은 순서였다.}

\begin{quote}
\textbf{조항} --- \textbf{이미 채택된 것을 되돌리려 할 때는, 되돌리는
근거가 기대는 사실을 코드나 자료에서 먼저 확인한다.} 기억으로 적은
전제 위에서 되돌리지 않는다.
\end{quote}

\section*{남은 것}

\begin{itemize}
\item \textbf{어느 범위가 배포에 맞나} --- 노트 468 이 물은 그대로 남았다.
사람에게 물을 것이다.
\item \textbf{판에도 갈래 안 점수를 같이 찍어 볼 만하다} --- 지금은
도메인 전체 하나뿐이라, 도메인 안에서 갈래 사이 정보에 얼마나 기대는지
안 보인다.
\item \texttt{OOC\_DROP} 은 계속 미결(노트 463) --- 이번 일로 바뀐 것은
없다.
\end{itemize}

\end{document}
